% 第 13 章　自然会“选择最省事的路”吗
\chapter{自然会“选择最省事的路”吗}

\step{反常识开场}

\begin{mainline}
把一根筷子斜插进一杯水里，从侧面看，筷子在水面处“折”了——水下那一截好像被谁往上掰了一下。没人碰它，杯子也没动，可你就是亲眼看见一根直筷子变成了两段。把筷子抽出来，它还是直的；放回去，又“折”了。

再来看一个海滩上的场面。救生员发现浅水区有人呼救，她离呼救者之间隔着一段干沙滩和一段海水。她在沙滩上能跑得飞快，可一下水，速度立刻慢下来。你猜她会怎么做？几乎没有一个有经验的救生员会沿着直线冲过去——那条路看起来最短，可它在慢吞吞的水里占的长度太长。她会先在沙滩上斜着多跑一段，再拐个弯下水。她没有纸笔，没有计算器，但她选的路线恰好是用时最短的路线。

现在，真正让人坐不住的事情来了：光从空气射进水里，也会在水面拐一个弯——而它拐弯的方式，和救生员拐弯的方式，竟然是同一种。好像光也“知道”自己在空气里快、在水里慢，所以宁愿在空气里多走一点，少在水里泡一会儿，最后走一条费时最少的路。可光没有眼睛，没有大脑，它出发的时候连终点的情况都还没“看到”，它凭什么“选”出了最省时间的路？

还不止光。镜子反射时，入射角永远等于反射角，像精心设计过的；一根晾衣绳松垮垮挂在那里，它的形状看似随意，却精确得可以用一个“最省”的原则算出来；你把篮球抛向空中，它走的那条抛物线，也藏着同样的秘密。这一章要问的就是：自然真的会“选择最省事的路”吗？如果会，它用的“省”字是什么意思？如果不会，那这些“巧合”又是怎么回事？
\end{mainline}

\step{先猜一猜}

老规矩，先把你的直觉押在桌上，别偷看后面。

第一题。一束光从空气斜射进水里，它在水中会朝哪个方向偏？（a）更靠近水面，几乎贴着水面走；（b）更靠近竖直方向，像被水面“吸”得直了一点；（c）不偏，方向不变。

第二题。沙滩上救生员在甲处，呼救者在水中乙处。岸边是一条直线。三条候选路线：甲到乙的直线；先在岸上多跑一段、再几乎垂直下水；先在岸上少跑、早早下水斜着游。哪条最快？或者你需要知道什么信息才能判断？

第三题。把一根两端固定的绳子松垮垮挂着（晾衣绳），它的形状最接近什么？（a）一段圆弧；（b）一段抛物线；（c）都不是。如果“自然喜欢最省”，你觉得绳子“省”的是什么——长度？高度？还是别的？

写下三个答案和理由。本章结尾我们逐一对账。提醒一句：这三题表面上是光学、营救和绳子，互不相干；如果你觉得它们背后可能有同一个机关在转动，请把你的猜测也写下来——那正是本章要挖的东西。

\step{动手实验}

\begin{experiment}[激光笔射进水里：量一量光拐的弯]
\textbf{材料}：一支激光笔（或强光小手电）、一只透明直壁玻璃杯、水、一滴牛奶（没有也行）、一张白纸、一支笔、一个量角器。

\textbf{步骤一}：杯里装大半杯水，滴入一滴牛奶搅匀（让光在水里显出路径；没有牛奶，就在杯后竖一张白纸观察光点）。把杯子放在白纸上，关灯，让激光笔贴着纸面从空气斜射向杯壁。你会清楚看到光在进入水的一瞬“折”了一下：水中的那段更靠近竖直方向。

\textbf{步骤二}：在纸上描下光在空气中和水中的两段路线，延长到入射点，量出两条路线与界面垂线（法线）的夹角，分别记作 \(\theta_1\)（空气中）和 \(\theta_2\)（水中）。改变入射的倾斜程度，再测三四组。你会发现：\(\theta_2\) 总是小于 \(\theta_1\)，而且两个角的正弦之比 \(\sin\theta_1/\sin\theta_2\) 几乎是个定数，大约 \(1.3\) 上下。

\textbf{步骤三}：把激光笔贴着杯底从水向空气射（杯子举高，从下方斜射）。光的弯折方向反了过来：进入空气后远离法线。把入射角逐渐加大，到某个角度后光突然不再射出，而是全部折回水里——像水面变成了一面镜子。

\textbf{步骤四}（加分项）：把水杯换成装水的方形塑料盒，让光先后穿过“空气—水—空气”两个平行界面。观察出射光与最初的入射光：方向完全一样，只是整体平移了一段。两块平行的玻璃不改变光的方向——这就是为什么隔着平整的窗户看外面，世界不会变形。

\textbf{记录}：把几组 \(\theta_1\)、\(\theta_2\) 记下来。本章后面要回答：这个约等于 \(1.3\) 的定数是什么？它和“最省时间”有什么关系？最后那个“水面变镜子”的现象，又凭什么从同一条规律里长出来？
\end{experiment}

还有一个零成本观察：晚上走过亮着路灯的积水路面，或者白天看一杯水里的吸管，随时留意“水下的东西看起来比实际浅”。把一个硬币丢进不透明杯子，退到刚好看不见它的位置，然后往杯里倒水——硬币“浮”进了视野。先记住这份惊讶。

\step{旧模型遇到麻烦}

我们先有的模型是“光走直线”。它在空气里、在水里各自都成立——你在实验里看到的两段确实都是直的。麻烦出在交界面上：为什么偏偏在那里拐弯？朝那个方向拐？“光走直线”对这些问题一个字也答不上来。

那就升级一下。两千年前，亚历山大港的希罗就注意到一件妙事：镜子反射时入射角等于反射角，这恰好等价于“光走的是从出发点经镜面到终点的所有路线里\emph{最短}的一条”。（把终点翻到镜子另一侧变成“镜像”，直线距离最短，一展开正好给出两角相等。）于是旧模型打上一个漂亮的补丁：\emph{自然走最短的路}。

可惜补丁经不起水。折射根本不服从“最短”：从空气到水里某点，最短的路是直线，光却偏不直走。力学这边也一样拆台：竖直上抛一个球，它明明可以“走直线最快到顶”，可它偏偏先快后慢、画一条抛物线；一只篮球飞向篮筐，走的不是直线而是弧线。“最短”这个标准，在折射和抛体面前双双失效。

更要命的是，每次打补丁都要新发明一个“自然在乎的东西”：反射在乎路程最短，折射在乎别的东西，抛体又在乎第三个东西……这不是一个理论，这是一堆事后诸葛亮。还有一个更根本的别扭：无论是“直线”还是“最短”，描述的都是\emph{空间的形状}，可抛球问题里起关键作用的明明是\emph{时间}——球在高点待得久、在低点跑得快，这些节奏信息被“最短”两个字完全丢掉了。我们需要一个统一的标准，而且必须先搞清楚一件事：自然“省”的到底是什么。

\step{发明新概念}

\begin{mainline}
\textbf{第一个新概念：光在不同介质里有不同的速度。}这是实验事实，不是假设：光在真空里约每秒三十万千米，在水里约每秒二十二万五千千米——只有真空中的四分之三左右；在玻璃里更慢。一旦接受这件事，“最短”和“最快”就分家了：救生员的沙滩上跑得快、水里游得慢，所以值得在快的地方多走几步，把慢的那段缩到最短。光也一样——如果光“在乎”的是时间，那么它在空气里多绕一点、少在水里泡一点，就完全是理性的选择。费马在十七世纪把这个想法立成原理：\emph{光走的路线，是所有可能路线中费时取稳定值的那条}。这叫费马原理。

\textbf{第二个新概念：把“一整条路径”当成一个整体来比较。}这是本章真正的思维升级。在此之前，力学是这样看世界的：盯着物体“此刻”的位置、速度和受力，算出“下一刻”往哪走，一步一步推着走。而现在，我们把整条从出发到到达的假想路径摆上桌，给每一条路径打一个总分，然后比较哪个分数最“稳”。前者像开车时只看前方三米，后者像出发前摊开整张地图。令人惊异的是，这两种看法最后给出完全相同的运动——它们是同一套自然规律的两种语言，而这一章学的是“地图语言”。

\textbf{第三个新概念：稳定值不等于最小值。}这是本章最容易被“最省事”三个字带偏的地方。想象山间的地形：盆地的底部是最低点，往任何方向挪一步都会升高——这是“最小”。但还有一种地方：两个山头之间的山口（鞍点）。站在山口正中央，沿山脊方向挪一点会升高，沿山谷方向挪一点会降低，而无论往哪边挪，高度的一阶变化都是零——山口也是一个“稳”的地方。费马原理说的“稳定值”就是这种意思：\emph{把真实路径轻轻挪动一点，总费时的一阶变化为零}。它像什么？像过山口时那一刻的“脚下平”：你在山口中央，感觉不到任何方向的坡度。它不像什么？它不像“考试分数越低越好”的排名赛——自然并不追求“最少”，它只要求“稳”。事实上，光走凹面镜反射的某些路线时，取的竟是费时\emph{最长}的路径——这就是本章的反直觉反例，后面详谈。

\textbf{第四个新概念：作用量。}费马原理只管光。抛出去的篮球、绕行的行星、振荡的弹簧，有没有类似的“总分”？有。十八世纪的物理学家们发现：对力学系统，也可以给每条假想路径打一个分数，这个分数叫\emph{作用量}；真实的运动让作用量取稳定值。这里必须立刻泼两盆冷水。第一盆：作用量\emph{不是}能量，“作用量取稳定”绝不等于“能量最小”——晾衣绳那才是能量（势能）最小的形状，两者是完全不同的原理。第二盆：它也不等于日常说的“省力”“省事”——有时真实路径的作用量比邻近路径大，有时小，有时是个山口式的鞍点，唯一不变的是“一阶不动”。所以本章标题里的“最省事”，严格的版本应该叫“最稳事”。

作用量的直觉可以先赊账感受一下（公式在数学桥层）。这个分数大致是“动能减势能”在整段旅程上的累加。抛一个球：走直线直奔最高点，意味着在势能最高的地方慢慢磨，动能被压得太低，“动能减势能”在那里难看得要命；贴着地面低飞，势能是省了，可要在同样的时间里赶到终点，就得飞得飞快，动能又花冒了。真实的抛物线是两头的妥协：该快的时候快，该高的时候高，把整段旅程的“动能减势能”总账调到一个挪不动的位置。这不是球在“权衡”，而是说——只有这条路线，牛顿定律在每一瞬间都成立。
\end{mainline}

\begin{center}
\begin{tikzpicture}[scale=1.0]
  % 真实路径与扰动路径
  \fill (0,0) circle (2pt) node[left=3pt] {出发点};
  \fill (8,2.2) circle (2pt) node[right=3pt] {终点};
  \draw[very thick,blue!70!black] (0,0) .. controls (2,1.6) and (5,2.4) .. (8,2.2)
    node[midway,above=6pt] {真实路径（作用量稳定）};
  \draw[thick,dashed,red!70!black] (0,0) .. controls (2,0.2) and (5,0.6) .. (8,2.2)
    node[midway,below=8pt] {挪歪一点的假想路径};
  \draw[-{Stealth[length=2.5mm]},gray!70!black] (3.1,1.85) -- (3.1,0.75);
  \node[gray!70!black,right=2pt] at (3.3,1.3) {轻轻挪动};
\end{tikzpicture}
\end{center}

变分思想就在这“轻轻挪动”四个字里：取真实路径，把它向任意方向挪歪一点点（出发点和终点不许动），算总分变了多少。若总分的一阶变化恒为零，这条路就是自然选中的路。“一阶”是关键词：挪一毫米，总分的变化不是零，但跟毫米\emph{成正比}的那部分是零，只剩下跟毫米的平方、立方成正比的残渣——就像站在山口中央，挪一小步的高度差随步长平方增长，步长越小越显得平。数学桥层会把这句话写成式子。

补一段极简历史，免得你以为这套想法是某天从天上掉下来的。十七世纪费马为光立下“最省时间”原理；十八世纪，莫佩尔蒂猜想着力的世界里也存在一个“最省”的量，并给它起名“作用量”，欧拉随即给出了严格的数学说法；再往后，拉格朗日和哈密顿把它锤炼成两套完整的力学语言——那正是第 14、15 章的主角。名字一路沿用下来，所以你在任何一本教科书里读到的“最小作用量原理”，其实都是本章这位“最稳”的老朋友。

\step{一句话模型}

\begin{oneline}
自然不“偷懒”，也不“选择”；真实的运动是这样一条路：把它轻轻挪动一点，代价（光的时间，力学的作用量）在一阶近似下不变——是“最稳”，不是“最小”。
\end{oneline}

\step{数学翻译器}

\begin{mainline}
先把救生员的选择算出来，用真实的数量级。设岸边是一条直线，救生员在岸上甲处，呼救者在水中，离岸垂直距离 \(30\) 米，横向相距 \(40\) 米。

\begin{center}
\begin{tikzpicture}[scale=0.09]
  % 岸：上方为沙滩，下方为水
  \fill[yellow!20] (0,0) rectangle (50,18);
  \fill[blue!12] (0,0) rectangle (50,-32);
  \draw[thick] (0,0) -- (50,0) node[midway,above=2pt] {};
  \node at (46,14) {沙滩};
  \node at (46,-27) {水};
  % 甲（救生员，在岸边）与乙（呼救者）
  \fill (4,0) circle (0.7) node[above=2pt] {甲};
  \fill (44,-30) circle (0.7) node[below=2pt] {乙};
  % 最优路径（沿岸跑约 31 米再下水）
  \draw[very thick,red!70!black] (4,0) -- (35,0);
  \draw[very thick,red!70!black] (35,0) -- (44,-30);
  % 直线路径
  \draw[thick,dashed,blue!70!black] (4,0) -- (44,-30);
  % 标注
  \node[red!70!black] at (20,2.5) {岸上多跑（快）};
  \node[red!70!black,rotate=-73] at (37.2,-15) {水里少泡（慢）};
  \node[blue!70!black,rotate=-53] at (24,-16) {直线：最短但不最快};
\end{tikzpicture}
\end{center}

救生员在沙滩上跑每秒 \(5\) 米，在水里游每秒 \(1.5\) 米。她在岸上跑到距甲 \(x\) 米处下水，再直线游向呼救者。总时间为
\[
  T(x)=\frac{x}{5}+\frac{\sqrt{30^2+(40-x)^2}}{1.5}\ \text{秒}.
\]
现在比较几条路线。沿直线冲（\(x=0\)，立刻下水）：岸上时间零，水中距离 \(\sqrt{30^2+40^2}=50\) 米，\(T=50/1.5\approx 33.3\) 秒。先在岸上跑满 \(40\) 米再垂直下水：\(T=8+30/1.5=28\) 秒——已经快了五秒多。再把 \(x\) 取细一点，逐个代入：\(x=25\) 米时 \(T\approx27.4\) 秒，\(x=27\) 时约 \(27.2\)，\(x=29\) 时约 \(27.1\)，\(x=31\) 时约 \(27.08\)，\(x=33\) 时约 \(27.1\)，\(x=35\) 时约 \(27.3\)。最低点在 \(x\approx31\) 米附近——比直线冲刺快约六秒。对溺水者来说，六秒就是生死。

请盯住这组数字的一个特征：在最优点附近，\(x\) 挪动整整两米，总时间只动零点零几秒；离最优点远了，同样的两米挪动带来的时间变化明显变大。这就是“一阶不变”在数字里的长相：\emph{谷底必然是平的}。反过来说，“找到那个怎么挪都不太变的位置”和“找到最低点”是同一件事——这就是变分思想的工作方式：不用真的算出最低点，只要找到“挪不动”的那个位置。请注意这笔账的结构：在快介质里多走约 \(31\) 米，只为把慢介质里的路程从 \(50\) 米压到约 \(31\) 米——\emph{用快路换慢路}，这就是折射拐弯的全部秘密。

把这个“用快换慢”的条件推到一般情形，就是折射定律。设光在两种介质中的速度分别为 \(v_1\)、\(v_2\)，入射角与折射角为 \(\theta_1\)、\(\theta_2\)（都相对法线测量），“总费时取稳定值”要求
\[
  \frac{\sin\theta_1}{\sin\theta_2}=\frac{v_1}{v_2}.
\]
真空中光速记为 \(c\)，把每种介质的“慢”程度写成 \(n=c/v\)（叫折射率），式子变成更对称的样子：
\[
  n_1\sin\theta_1=n_2\sin\theta_2.
\]
这就是斯涅尔定律。它回答四个问题。描述什么：光在两种介质界面上拐多少。为什么这样写：因为只有这样，总费时的一阶变化才为零——两边相等正是“挪一挪入射点不再省时间”的条件。何时成立：两种均匀介质的平整界面，且光的波长远小于界面尺度。何时失效：光被小到与波长相当的结构摆弄时（衍射光栅、肥皂泡的颜色），“一条路径”的图像整个垮掉，得用波动光学重做。

新公式过三道例行检查。第一，\(n_1=n_2\)（两边是同一种介质）：\(\theta_1=\theta_2\)，光不拐弯——退化成“光走直线”，正确。第二，\(n_2\) 极大（光在里面慢得像爬）：\(\sin\theta_2\) 被压得极小，光几乎沿法线走——慢的路段被缩到最短，和救生员“垂直下水”的极限一致。第三，方向反转：让光从水射回空气，\(\theta\) 的大小关系整个对调，公式照样成立——光路可逆，你在实验步骤三里亲眼见过。
\end{mainline}

还有那个“水面变镜子”：让光从水（\(n_1\approx1.33\)）射向空气（\(n_2\approx1\)），斯涅尔定律给出 \(\sin\theta_2=1.33\sin\theta_1\)。当 \(\theta_1\) 大到约 \(49^{\circ}\)，右边到达 \(1\)——\(\theta_2=90^{\circ}\)，折射光贴着水面溜走；再加大入射角，方程要求 \(\sin\theta_2>1\)，无解。数学的“无解”就是物理的“全反射”：光百分之百折回水中。潜水员抬头看，只有头顶一个亮圆窗，窗外的水面全是镜子，人称“斯涅尔窗”。光纤能把光锁在细玻璃丝里拐着弯送出去几千米，靠的也是同一件事。

顺带做一次带真实数据的估算，看看“光在介质里变慢”在工程里有多实在。光纤里光的等效速度约每秒二十万千米。从北京到上海，光纤长度按一千三百千米计，单程延迟约 \(1300/200000\) 秒，即 \(6.5\) 毫秒，往返约 \(13\) 毫秒——这就是跨省网游延迟里“怎么也砍不掉”的那一部分：交换机再快、软件再优化，光在玻璃里的速度替延迟画了一条硬地板。同一个原理换个舞台：隔着一块三毫米厚的窗玻璃，光比走真空多花约 \((1.5-1)\times 0.003/(3\times10^{8})\) 秒，即五皮秒（一皮秒是一万亿分之一秒）。五皮秒短到荒谬，可精密光学实验和原子钟授时系统偏偏要把这种账一笔笔记清楚。

\begin{mathbridge}
力学版本的总分这样写。对质量 \(m\) 的物体，动能记 \(T=\frac12 mv^2\)，势能记 \(V\)（两者第 9 章已有），每条假想路径 \(x(t)\) 的作用量定义为
\[
  S[\,x(t)\,]=\int_{t_1}^{t_2}\bigl(T-V\bigr)\,dt,
\]
即在整段旅程上把“动能减势能”逐刻累加起来。真实路径满足
\[
  \delta S=0,
\]
读作：把路径换成 \(x(t)+\varepsilon\,\eta(t)\)（\(\eta\) 是任意扰动，两端为零，\(\varepsilon\) 是无穷小量），\(S\) 的变化在一阶 \(\varepsilon\) 上为零。注意 \(T-V\) 的减号：这不是总能量 \(T+V\)，所以“作用量稳定”从头到尾都不是“能量最小”。

这个条件为什么等价于牛顿定律？直觉是：在某一小段时间里把路径向上挪一点，势能 \(V\) 那段的积分立刻跟着变（位置变了），动能 \(T\) 那段的积分则通过速度的变化跟着变；要让总账的一阶变化为零，两件事必须精确抵消——逐点写下来，抵消条件正是 \(F=ma\)。所以最小作用量原理不是牛顿定律之外的“第二条自然规律”，它是同一枚硬币的另一面：牛顿定律是逐时刻的“局部语言”，作用量原理是整段旅程的“整体语言”，二者互相翻译，一字不差。

再看一眼光的版本有多像：费马原理里被累加的量是逐小段的“走时” \(n\,ds/c\)，作用量原理里被累加的是逐时刻的 \(T-V\)。两章之后（第 15 章）你会发现，连“被累加什么”都可以统一到一个更一般的框架里。
\end{mathbridge}

\begin{challenge}
“稳定”而非“最小”，在数学上就是“一阶变分为零”。把 \(S[x+\varepsilon\eta]\) 对 \(\varepsilon\) 展开，令一阶项对任意 \(\eta(t)\) 为零，可以推出一条逐点成立的微分方程——欧拉方程，第 14 章会沿着这条路把拉格朗日力学整个建起来。

这里只埋一颗通向远方的种子。费曼后来给“光怎么知道哪条路最省时”一个量子力学的回答：光（以及一切微观客体）根本不“知道”，它也并不只走一条路——在量子描述里，从出发到到达要计入\emph{所有}路径，每条路径贡献一个带相位的振幅，相位由作用量决定。在稳定路径附近，邻近路径的作用量几乎相同，相位几乎一致，振幅互相加强；远离稳定路径，相位乱成一团，彼此抵消。于是宏观世界里“光选了最稳的路”，是亿万条路径集体投票后的幸存者。第 42 章之后，这个故事会正式展开；今天只需记住：本章的“稳定值条件”将来会被证明比牛顿定律活得更久——它在量子世界里依然成立，而牛顿定律不会。
\end{challenge}

\step{模型的边界}

\begin{mainline}
先把适用范围钉在墙上，再清点误用。

费马原理的适用条件：介质均匀或缓慢变化（渐变介质里光走连续曲线，海市蜃楼就是例子）；界面尺度远大于波长。失效方向：波长尺度的结构（衍射）、光强极强时的非线性介质，都要换理论。

作用量原理的适用条件：系统的力都能写成势能（重力、弹力、静电力可以），否则标准框架装不下。摩擦力这类耗散力就是个麻烦户：搓热的双手损失的能量跑去哪里，作用量原理自己说不清楚，得请热力学帮忙。它不是错了，而是管不了“耗散”这本账。还有一条隐性条件：出发点和终点（含时刻）必须先固定，比较才有意义——“从家到公司哪条路最稳”是合法问题，“随便去哪儿哪条路最稳”则根本不是问题。
\end{mainline}

典型误用清单：

\begin{itemize}[leftmargin=1.8em]
  \item \emph{“自然有目的，它想偷懒。”}这是拟人化的解释，不是实验事实。实验事实是：稳定值条件与运动方程等价，多一个字都是文学。正确的因果方向也恰好相反：不是自然先“看完全局再选路”，而是每个瞬间都服从局部规律，无数局部叠加起来，恰好让整条路径满足稳定条件。就像蚁群没有谁规划全局，却走出了高效路线——宏观的“智慧”是微观的死规矩堆出来的。
  \item \emph{“作用量最小等于能量最小。”}完全两回事。被累加的是 \(T-V\)，不是 \(T+V\)；真实路径的作用量甚至常常比邻近路径\emph{大}。要能量最小请去找晾衣绳（那是静力平衡，势能最小），别来找作用量。
  \item \emph{“最短时间原理。”}名字本身就错了一半。凹面镜成像的某些光路取的是时间\emph{最长}的路径，平面镜某些绕远的合法光路是鞍点。严格的说法永远是“稳定值”。凡是你读到“最小作用量原理”的地方，心里自动替换成“稳定作用量原理”，就不会被骗。
  \item \emph{“既然整体语言等价于局部语言，那以后不用学牛顿定律了。”}不行。整体语言好算约束多的问题（第 14 章的主场），但受力分析、接触、碰撞的瞬间细节，仍要回到力和加速度。两种语言是互相备份，不是互相取代。
\end{itemize}

\step{回头解释开场现象}

现在对账。先答“先猜一猜”。

第一题：光从空气进水中，\(\theta_2<\theta_1\)，更靠近竖直方向，选（b）——你在实验里量的几组数据和那个约 \(1.3\) 的比值（水的折射率）已经替你确认了。

第二题：救生员路线取决于两个速度之比；岸上跑得越快、水里游得越慢，越该在岸上多绕。本章算过：每秒 \(5\) 米对每秒 \(1.5\) 米时，最优路线比直线快约六秒。这和光走的是同一笔账，只是“折射率”换成了速度比。

第三题：晾衣绳的形状既不是圆弧也不是抛物线，而是一条叫“悬链线”的曲线（只在绳子绷得较紧、垂度较小时近似抛物线）。它“省”的是重力势能：在两端固定、长度固定的所有假想形状里，真实挂着的形状让整根绳子的重心最低——任何一点被挪高或挪低，总势能都会升高。这是“稳定（最小）势能”而不是“最小作用量”，但思考方式完全同构：把整条曲线当整体，轻轻挪动，一阶不变。蜘蛛网、高压输电线、悬索桥的主缆，全是这条曲线的工程分身：桥梁设计师把桥面载荷挂上去，让缆索自己找到那个“稳”的形状，材料就只需承受拉力，不用硬扛弯矩。

开场里还漏了一面镜子没对账：反射时入射角为什么等于反射角？用费马原理一句话还清——在均匀介质里速度不变，省时就是省路，而从出发点经镜面到终点的最短路线，恰好就在两角相等处（挑战题第一题请你亲手证一遍）。至此，反射、折射、抛体、绳子，四种风马牛不相及的现象，被同一句“挪一挪，一阶不变”串在了一根线上。

筷子“折”了：光从筷子上出发，穿过水面进入空气时按斯涅尔定律远离法线拐了个弯；你的眼睛只会沿直线往回追，于是把水下那截“看”到了更浅、更靠上的位置。泳池看起来比实际浅约四分之一，退到杯外看不见的硬币加水后“浮”出来，都是同一笔账。海市蜃楼则是渐变版：贴近热路面的空气折射率略小，光在一层层空气里连续地“用快换慢”，走出一条向上弯的弧线，把天空的光送到你眼睛里，看起来像一汪水。

还有一类更安静的应用藏在你脸上和口袋里：眼镜、相机镜头、望远镜。镜片设计者的全部工作，就是用不同形状、不同折射率的玻璃，把光线一段一段的“走时”重新分配，让从同一个物点出发、穿过镜片不同位置的光，到达像点时费时全都相等——各条光路“一样稳”，才能齐刷刷会聚成一个清晰的点；哪一片玻璃的曲率算错了，某条光路就提前或拖后，像就糊了。你手机上那颗小小的镜头里有五六片形状古怪的镜片，每一片都在做这道“稳定值”作业。

至于本章标题的问题——自然会“选择最省事的路”吗？答案是：不会“选择”，也没有什么“省事”；有的只是一条对“轻轻挪动”无动于衷的路径。开场那三种“巧合”——反射的两角相等、折射的定数、抛体的弧线——从这一句里全部长了出来。

\step{脑洞实验与挑战题}

\begin{thought}[把太阳遮住一半：光在极端尺度上还走“最稳的路”吗]
日全食时，星光擦着太阳边缘射向地球。按牛顿的直觉，光该走直线；可实际观测到星光被太阳引力掰弯了一个微小但确定的角度。广义相对论（第 41 章）给出的读法和本章惊人地同构：光仍然走“最稳的路”，只不过在太阳附近，时空本身被质量压弯了，“最稳”的标准要用弯曲时空里的尺子来量——那条路叫测地线，是弯曲世界里“尽可能直”的路。换句话说，本章的“地图语言”在极端尺度上不但没死，反而升级成了引力理论本身：行星绕日，也是沿着弯曲时空里的测地线“自由滑行”，根本不需要谁在中间拉一根线。一个关于“光为什么拐弯”的问题，最后竟把引力也收编了进来。

其实你在民航客机上已经见过“弯曲世界里的最稳路”。北京飞纽约的航线在平面地图上弯向北方、擦过北极附近，看起来绕了远；可地球是球面，地图是压平的，球面上“尽可能直”的路线投到平面地图上就是这副弯腰的样子。航空公司选它，不为看风景，只为它最短、最省油。球面上的测地线、太阳附近的星光、晾衣绳的悬链线，都是同一个句型：先把“稳”的尺子定义对，剩下的交给“轻轻挪动”。
\end{thought}

\begin{thought}[如果光真的同时走所有路]
把挑战层那颗种子泡一泡水。设想在光源和屏幕之间挡一块开了一条缝的板，缝越开越小，按费马原理光路还是那一条；可实验上缝小到波长量级时，光在屏幕上摊成一片衍射花纹——“一条路”的图像破产了。量子的回答是：光一直在走所有的路，只是宏观尺度下，稳定路径之外的路彼此抵消得一丝不剩；当你把缝开得很小，相当于人为删掉了大部分路径，剩下的路径不再抵消干净，稳定路径的“霸权”就露了馅。这个脑洞的意义在于：本章的原理在微观世界不是被推翻，而是被“扶正”——作用量从一条算账捷径，变成了量子理论的入场券（第 42 章起）。
\end{thought}

\begin{puzzles}
\item 用费马原理推出反射定律：光从甲点经平面镜反射到乙点，证明费时最短（此时也是路程最短）的路线恰好满足入射角等于反射角。\emph{思路提示}：把乙点关于镜面翻转到镜子背后，得到镜像点乙$'$。从甲经镜面任一点到乙的路程，等于从甲经同一点到乙$'$的路程。甲到乙$'$的最短路是什么？直线。这条直线与镜面的交点处，两个角什么关系？
\item 潜水员在水中抬头，看到的水面为什么只有一个圆形亮窗，窗外全是黑亮的“镜子”？估算这个圆窗的半顶角。\emph{思路提示}：用斯涅尔定律 \(n_1\sin\theta_1=n_2\sin\theta_2\)，从水到空气，令折射角 \(\theta_2=90^{\circ}\)，代入 \(n_1\approx1.33\)、\(n_2\approx1\)，解出临界角约为 \(49^{\circ}\)。比这个角度更斜的光全部全反射，到不了眼睛。
\item 游乐园里两条滑梯从同一高度出发、到达同一低点：一条是笔直的斜坡，一条是先陡降、凹下去再平缓收尾的弯道。哪条先到终点？直觉说“直线最短所以最快”，对吗？\emph{思路提示}：最短不等于最快——这正是本章的核心。弯道先把高度换成速度，全程平均速度更高，即使路程更长也可能先到（极限版叫“最速降线”）。再想想：如果弯道凹得太深、太平，会不会反而慢？“稳”的平衡点在哪里？
\item 一位游客站在湖边看对岸的树，既看到树的正影，又看到水面里的倒影。树顶发出的光，一条路径直奔眼睛，一条经水面反射到眼睛。两条路径都真实存在——费马原理不是“只走一条路”吗？\emph{思路提示}：稳定值可以不止一个，就像山口可以有好几个。直接路径和反射路径各自都是“挪一挪、一阶不变”的路，自然对每条都照走不误。这也解释了为什么镜面能成像：一束光里的每条光线都在解自己的“稳定值题”。
\end{puzzles}

至此，我们手里多了一种全新的语言：不再逐时刻追问“下一瞬往哪走”，而是摊开整条路径问“哪条路最稳”。这条语言在折射、反射、抛体、晾衣绳身上全都灵验，可它真正的威力还没有释放——下一章我们将看到，只要把“被累加的量”写成 \(T-V\)，再让路径轻轻挪动，\(F=ma\) 会自己从“稳定值条件”里长出来，而且许多让牛顿语言焦头烂额的复杂系统（转动的、被约束的、多自由度的），在新语言里会变得出奇地简单。那就是拉格朗日力学。
